\documentclass[10pt,a4paper]{amsart}
\usepackage[margin=19mm]{geometry}
\usepackage{amsmath,amssymb,mathtools}
\usepackage[scr=rsfs,cal=euler]{mathalfa}
\usepackage{xfrac}
\usepackage[unicode,hidelinks,bookmarks=false]{hyperref}
\newcommand{\ie}{i.e.\ }
\newcommand{\cf}{cf.\ }
\newcommand{\resp}{respectively\ }
\newcommand{\Subsection}{\subsection}
\providecommand{\nobreakdash}{\nobreak\hbox{-}}
\setlength{\emergencystretch}{2em}
\multlinegap0pt
\newtheorem{lemma}{Lemma}
\newtheorem{theorem}{Theorem}
\title{Constitutive Modelling of Korteweg Fluids Using Liu's Method}
\author{Zagorka Mati\'c, Srboljub Simi\'c, Peter V\'an}
\date{Source: arXiv:2604.26846v1 (CC BY 4.0)}
\begin{document}
\maketitle

\noindent\emph{Source and licence.} Constitutive Modelling of Korteweg Fluids Using Liu's Method. Version arXiv:2604.26846v1 (CC BY 4.0), distributed by arXiv under the Creative Commons Attribution 4.0 International licence, http://creativecommons.org/licenses/by/4.0/, as recorded in the arXiv licence metadata for this exact version and shown on the abstract page https://arxiv.org/abs/2604.26846v1. Full licence notice: \texttt{licenses/arxiv-2604.26846-CC-BY-4.0.txt}.\par\smallskip

\noindent\textit{Editorial scope.} Counted unit: the article's theorem fluidi (source0.tex lines 1155--1169). The article's complete original proof of it is delivered with it (source0.tex lines 1171--1282, one \texttt{proof} environment of 112 lines). No mathematics is written, repaired or re-derived: the statement and the proof are the article's own lines, in the article's own notation, and no retained line is substituted or re-flowed.  Retained uncounted as the source-local context the proof needs: lines 881--897; lines 923--925; lines 1030--1039; lines 1076--1093.  Nothing was omitted from any retained span, so the package carries no omission record.  No retained line contains a citation, so the wrapper carries no bibliography and no bibliography entry.  No retained line contains a figure environment, an image-inclusion command or drawing code, so no image is delivered and the package declares no asset dependency.  Editorial formatting only, in addition: an AMS \texttt{amsart} preamble that restates the article's own declarations and macros used by the retained lines, namely the environments \texttt{lemma}, \texttt{theorem} on separate counters, together with the standard packages those lines need.  The retained environments print in this document as Lemma 1, Theorem 1 in order of appearance, whereas the article prints the same items with the numbering of its own full document.  The complete native source of the article as distributed in the arXiv e-print for this exact version is bundled unchanged as \texttt{sources/199426/source0.tex}.
\begin{align} \label{osnovna rezidualna nejednakost Kortevegovi fluidi}
	& \dfrac{\partial\rho}{\partial x_j} \left[sv_j+\rho\overline{\dfrac{\partial s}{\partial\rho}}v_j+\dfrac{\partial\varphi_j}{\partial\rho}-\Lambda^\rho v_j-\Lambda^{v_i}v_iv_j+\Lambda^{v_i}\dfrac{\partial t_{ij}}{\partial\rho}
	\right. 
	\nonumber \\  
	& \qquad \left. -\dfrac{1}{2}\Lambda^e|\mathbf{v}|^2v_j -\Lambda^e e v_j-\Lambda^e \rho\dfrac{\partial e}{\partial\rho}v_j+\Lambda^e\dfrac{\partial t_{ij}}{\partial\rho}v_i-\Lambda^e\dfrac{\partial q_j}{\partial\rho}\right] 
	\nonumber \\
	& \quad + \dfrac{\partial\rho_{,k}}{\partial x_j} \left[\rho\overline{\dfrac{\partial s}{\partial\rho_{,k}}}v_j+\dfrac{\partial\varphi_j}{\partial\rho_{,k}}+\Lambda^{v_i}\dfrac{\partial t_{ij}}{\partial\rho_{,k}}-\Lambda^e\rho\dfrac{\partial e}{\partial\rho_{,k}}v_j % \right. 
	% \\ & \quad \left. 
	+ \Lambda^e\dfrac{\partial t_{ij}}{\partial\rho_{,k}}v_i-\Lambda^e\dfrac{\partial q_j}{\partial\rho_{,k}}-\Lambda^{\nabla\rho}_k v_j\right] 
	\\
	& \quad +\dfrac{\partial\theta}{\partial x_j}\left[\rho\overline{\dfrac{\partial s}{\partial\theta}}v_j+\dfrac{\partial\varphi_j}{\partial\theta}+\Lambda^{v_i}\dfrac{\partial t_{ij}}{\partial\theta}-\Lambda^e\rho\dfrac{\partial e}{\partial\theta}v_j+\Lambda^e\dfrac{\partial t_{ij}}{\partial\theta}v_i-\Lambda^e\dfrac{\partial q_j}{\partial\theta}\right] 
	\nonumber \\
	& \quad + D_{ij}\left[\rho s \delta_{ij}-\Lambda^\rho\rho\delta_{ij}-\Lambda^{v_i}\rho v_j-\Lambda^{v_k}\rho v_k\delta_{ij}-\Lambda^e\rho v_i v_j-\dfrac{1}{2}\rho\vert \mathbf{v}\vert^2\Lambda^e\delta_{ij} \right.
	\nonumber \\ 
	& \qquad \left. - \Lambda^e\rho e\delta_{ij} + \Lambda^et_{ij}-\Lambda^{\nabla\rho}_k\dfrac{\partial\rho}{\partial x_k}\delta_{ij}-\Lambda^{\nabla\rho}_j\dfrac{\partial\rho}{\partial x_i} \right] = \Sigma\geq 0. 
	\nonumber 
\end{align}

	\begin{equation} \label{mnozitelj Lambda-e Kortevegovi fluidi}
		\Lambda^{e} = \dfrac{1}{\theta}. 
	\end{equation}

	\begin{align}
		\label{X_jrho_,kl prim Kortevegovi fluidi}
		& \dfrac{\partial\varphi_j}{\partial\rho_{,kl}}-\Lambda^e\dfrac{\partial q_j}{\partial\rho_{,kl}}=0, 
		\\
		\label{X_jtheta_,k prim Kortevegovi fluidi}
		& \dfrac{\partial\varphi_j}{\partial\theta_{,k}}-\Lambda^e\dfrac{\partial q_j}{\partial\theta_{,k}}=0, 
		\\
		\label{X_jD_kl prim Kortevegovi fluidi}
		& \dfrac{\partial\varphi_j}{\partial D_{kl}}-\Lambda^e\dfrac{\partial q_j}{\partial D_{kl}}-\Lambda^{\nabla\rho}_j\rho\delta_{kl}=0.
	\end{align}

\begin{lemma} \label{Lema:Rezidualna nejednakost 2}
	Residual inequality \eqref{osnovna rezidualna nejednakost Kortevegovi fluidi} reduces to
	\begin{align} \label{rezidualna nejednakost 2 Kortevegovi fluidi}
		& \dfrac{\partial\rho}{\partial x_j} \left[\dfrac{\partial\varphi_j}{\partial\rho} 
		- \Lambda^e\dfrac{\partial q_j}{\partial\rho}\right] 
		+ \dfrac{\partial\rho_{,k}}{\partial x_j} \left[\dfrac{\partial\varphi_j}{\partial\rho_{,k}} 
		- \Lambda^e\dfrac{\partial q_j}{\partial\rho_{,k}} \right] 
		%\nonumber \\
		+ \dfrac{\partial\theta}{\partial x _j} \left[ \dfrac{\partial\varphi_j}{\partial\theta} 
		- \Lambda^e \dfrac{\partial q_j}{\partial\theta} \right]  
		\\
		& \quad + D_{ij} \left[ -\rho^2\delta_{ij} \left(\overline{\dfrac{\partial s}{\partial\rho}} 
		- \Lambda^e \dfrac{\partial e}{\partial\rho}\right) + \Lambda^et_{ij} 
		- \Lambda^{\nabla\rho}_k\dfrac{\partial\rho}{\partial x_k}\delta_{ij} 
		- \Lambda^{\nabla\rho}_j\dfrac{\partial\rho}{\partial x_i} \right] = \Sigma \geq 0. 
		\nonumber 
	\end{align}
\end{lemma}

\begin{theorem}
	The final form of the residual inequality reads 
	\begin{align} \label{rezidualna nejednakost final Kortevegovi fluidi}
		& \dfrac{\partial}{\partial x_j} \left( \varphi_j - \Lambda^eq_j 
		- \Lambda^{\nabla\rho}_j \rho D_{kl}\delta_{kl}\right) 
		+ \dfrac{\partial\Lambda^e}{\partial\theta} \dfrac{\partial\theta}{\partial x_j}q_j 
		\\
		& \quad + D_{ij} \left[ \left( -\rho^2 \left( \overline{\dfrac{\partial s}{\partial\rho}} 
		- \Lambda^e \dfrac{\partial e}{\partial\rho} \right) 
		+ \dfrac{\partial\Lambda^{\nabla\rho}_k}{\partial x_k}\rho\right) \delta_{ij} 
		+ \Lambda^et_{ij} - \Lambda^{\nabla\rho}_j\dfrac{\partial\rho}{\partial x_i} \right]
		= \Sigma \geq 0. 
		\nonumber 
	\end{align}
\end{theorem}

\begin{proof}
	We shall start by transforming the first two terms in the residual inequality \eqref{rezidualna nejednakost 2 Kortevegovi fluidi}. Taking into account that entropy flux $\varphi_{j}$ and heat flux $q_{j}$ are determined by the constitutive functions
	\begin{align*}
		\varphi_j & = \varphi_{j}(\rho, \rho_{,k}, \rho_{,kl}, \theta, \theta_{,k}, D_{kl}), 
		\\
		q_j & = q_{j}(\rho, \rho_{,k}, \rho_{,kl}, \theta, \theta_{,k}, D_{kl}), \quad 
		j = 1, 2, 3. 
	\end{align*} 
	It follows 
	\begin{align*}
		\dfrac{\partial\varphi_j}{\partial x_j}
		& = \dfrac{\partial\varphi_j}{\partial\rho}\dfrac{\partial\rho}{\partial x_j} 
		+ \dfrac{\partial\varphi_j}{\partial\rho_{,k}}\dfrac{\partial\rho_{,k}}{\partial x_j} 
		\\
		& \quad + \dfrac{\partial\varphi_j}{\partial\rho_{,kl}}\dfrac{\partial\rho_{,kl}}{\partial x_j}+\dfrac{\partial\varphi_j}{\partial\theta}\dfrac{\partial\theta}{\partial x_j}+\dfrac{\partial\varphi_j}{\partial\theta_{,k}}\dfrac{\partial\theta_{,k}}{\partial x_j}+\dfrac{\partial\varphi_j}{\partial D_{kl}}\dfrac{\partial D_{kl}}{\partial x_j},
	\end{align*}
	and analogously for the heat flux. Taking this into account, first line in the residual inequality \eqref{rezidualna nejednakost 2 Kortevegovi fluidi} in Lemma \ref{Lema:Rezidualna nejednakost 2} can be transformed in the following way
	\begin{align*}
		& \dfrac{\partial\rho}{\partial x_j} \left[ \dfrac{\partial\varphi_j}{\partial\rho} 
		- \Lambda^e\dfrac{\partial q_j}{\partial\rho} \right] 
		+ \dfrac{\partial\rho_{,k}}{x_j} \left[ \dfrac{\partial\varphi_j}{\partial\rho_{,k}} 
		- \Lambda^e\dfrac{\partial q_j}{\partial\rho_{,k}} \right] \\
		& \quad = \dfrac{\partial\varphi_j}{\partial\rho} \dfrac{\partial\rho}{\partial x_j} 
		+ \dfrac{\partial\varphi_j}{\partial\rho_{,k}} \dfrac{\partial\rho_{,k}}{\partial x_j} 
		- \Lambda^e\left[\dfrac{\partial q_j}{\partial\rho}\dfrac{\partial\rho}{\partial x_j} 
		+ \dfrac{\partial q_j}{\partial\rho_{,k}} \dfrac{\partial\rho_{,k}}{\partial x_j}\right] \\
		& \quad = \dfrac{\partial\varphi_j}{\partial x_j} - \left(   
		\dfrac{\partial\varphi_j}{\partial\rho_{,kl}} \dfrac{\partial\rho_{,kl}}{\partial x_j} 
		+ \dfrac{\partial\varphi_j}{\partial\theta} \dfrac{\partial\theta}{\partial x_j} 
		+ \dfrac{\partial\varphi_j}{\partial\theta_{,k}} \dfrac{\partial\theta_{,k}}{\partial x_j} 
		+ \dfrac{\partial\varphi_j}{\partial D_{kl}} \dfrac{\partial D_{kl}}{\partial x_j} \right) \\
		& \qquad - \Lambda^e \left[ \dfrac{\partial q_j}{\partial x_j} 
		- \left( \dfrac{\partial q_j}{\partial\rho_{,kl}} \dfrac{\partial\rho_{,kl}}{\partial x_j} 
		+ \dfrac{\partial q_j}{\partial\theta}\dfrac{\partial\theta}{\partial x_j} 
		+ \dfrac{\partial q_j}{\partial\theta_{,k}} \dfrac{\partial\theta_{,k}}{\partial x_j} 
		+ \dfrac{\partial q_j}{\partial D_{kl}} \dfrac{\partial D_{kl}}{\partial x_j}\right) \right] \\
		& \quad = \dfrac{\partial\varphi_j}{\partial x_j} 
		- \Lambda^e \dfrac{\partial q_j}{\partial x_j} % \\ & \qquad 
		- \left[ \left( \dfrac{\partial\varphi_j}{\partial{\rho_{,kl}}} \dfrac{\partial\rho_{,kl}}{\partial x_j} 
		- \Lambda^e \dfrac{\partial q_j}{\partial{\rho_{,kl}}} \dfrac{\partial\rho_{,kl}}{\partial x_j} \right) 
		+ \left( \dfrac{\partial\varphi_j}{\partial{\theta}}\dfrac{\partial\theta}{\partial x_j} 
		- \Lambda^e\dfrac{\partial q_j}{\partial{\theta}} \dfrac{\partial\theta}{\partial x_j} \right) \right. \\
		& \qquad \left. + \left( \dfrac{\partial\varphi_j}{\partial{\theta_{,k}}} \dfrac{\partial\theta_{,k}}{\partial x_j} 
		- \Lambda^e \dfrac{\partial q_j}{\partial{\theta_{,k}}} 
		\dfrac{\partial\theta_{,k}}{\partial x_j}\right) 
		+ \left(\dfrac{\partial\varphi_j}{\partial D_{kl}}\dfrac{\partial D_{kl}}{\partial x_j} 
		- \Lambda^e \dfrac{\partial q_j}{\partial D_{kl}} 
		\dfrac{\partial D_{kl}}{\partial x_j} \right) \right]\\
		& \overset{\eqref{X_jrho_,kl prim Kortevegovi fluidi}, \eqref{X_jtheta_,k prim Kortevegovi fluidi}, \eqref{X_jD_kl prim Kortevegovi fluidi}} {=} 
		\dfrac{\partial\varphi_j}{\partial x_j} - \Lambda^e \dfrac{\partial q_j}{\partial x_j} 
		- \left( \dfrac{\partial\varphi_j}{\partial{\theta}}\dfrac{\partial\theta}{\partial x_j} 
		- \Lambda^e\dfrac{\partial q_j}{\partial{\theta}}\dfrac{\partial\theta}{\partial x_j} \right) 
		- \Lambda^{\nabla\rho}_j\rho\delta_{kl} \dfrac{\partial D_{kl}}{\partial x_j}.
	\end{align*}
	Using this transformation, inequality \eqref{rezidualna nejednakost 2 Kortevegovi fluidi} can be written in the form 
	%\begin{equation}
	\begin{align} \label{rezidualna nejednakost 3 Kortevegovi fluidi}
		& \dfrac{\partial\varphi_j}{\partial x_j} - \Lambda^e \dfrac{\partial q_j}{\partial x_j} - \Lambda^{\nabla\rho}_j\rho\delta_{kl}\dfrac{\partial D_{kl}}{\partial x_j}
		\\
		& \quad +D_{ij}\left[-\rho^2\delta_{ij}\left(\overline{\dfrac{\partial s}{\partial\rho}}-\Lambda^e\dfrac{\partial e}{\partial\rho}\right)+\Lambda^et_{ij}-\Lambda^{\nabla\rho}_k\dfrac{\partial\rho}{\partial x_k}\delta_{ij}-\Lambda^{\nabla\rho}_j\dfrac{\partial\rho}{\partial x_i}\right]=\Sigma\geq 0. 
		\nonumber 
	\end{align}
	%\end{equation}
	Having in mind that $\Lambda^e$ depends only on temperature $\theta$ (see equation \eqref{mnozitelj Lambda-e Kortevegovi fluidi}), residual inequality \eqref{rezidualna nejednakost 3 Kortevegovi fluidi} can be written as
	\begin{align*}
		& \dfrac{\partial}{\partial x_j}\varphi_j - \Lambda^e\dfrac{\partial}{\partial x_j}q_j 
			- \Lambda^{\nabla\rho}_j\rho\delta_{kl} \dfrac{\partial D_{kl}}{\partial x_j} 
			+ D_{ij} \left[ -\rho^2\delta_{ij}\left(\overline{\dfrac{\partial s}{\partial\rho}} 
			- \Lambda^e\dfrac{\partial e}{\partial\rho} \right) + \Lambda^et_{ij} 
			- \Lambda^{\nabla\rho}_k\dfrac{\partial\rho}{\partial x_k}\delta_{ij} 
			- \Lambda^{\nabla\rho}_j\dfrac{\partial\rho}{\partial x_i} \right] 
		\\
		& \quad = \dfrac{\partial}{\partial x_j} \left( \varphi_j - \Lambda^eq_j\right) 
			+ \dfrac{\partial\Lambda^e}{\partial x_j}q_j 
			- \Lambda^{\nabla\rho}_j\rho\delta_{kl} \dfrac{\partial D_{kl}}{\partial x_j} 
		\\ 
		& \qquad + D_{ij} \left[ \left( -\rho^2\left(\overline{\dfrac{\partial s}{\partial\rho}} 
			- \Lambda^e \dfrac{\partial e}{\partial\rho} \right) 
			- \Lambda^{\nabla\rho}_k\dfrac{\partial\rho}{\partial x_k} \right) \delta_{ij} 
			+ \Lambda^e t_{ij} - \Lambda^{\nabla\rho}_j\dfrac{\partial\rho}{\partial x_i} \right] \\
			& \quad = \dfrac{\partial}{\partial x_j} \left( \varphi_j - \Lambda^eq_j \right) 
			+ \dfrac{\partial\Lambda^e}{\partial x_j}q_j 
			- \dfrac{\partial}{\partial x_j} \left( \Lambda^{\nabla\rho}_j\rho D_{kl} \delta_{kl} \right) 
			+ \dfrac{\partial}{\partial x_j}\left(\Lambda^{\nabla\rho}_j\rho\right) \delta_{kl} D_{kl} \\
			& \qquad + D_{ij} \left[ \left( -\rho^2\left(\overline{\dfrac{\partial s}{\partial\rho}} 
			- \Lambda^e\dfrac{\partial e}{\partial\rho} \right) 
			- \Lambda^{\nabla\rho}_k\dfrac{\partial\rho}{\partial x_k} \right) \delta_{ij} 
			+ \Lambda^e t_{ij} - \Lambda^{\nabla\rho}_j\dfrac{\partial\rho}{\partial x_i} \right] \\
			& \quad = \dfrac{\partial}{\partial x_j} \left( \varphi_j - \Lambda^e q_j 
			- \Lambda^{\nabla\rho}_j \rho D_{kl}\delta_{kl} \right) 
			+ \dfrac{\partial\Lambda^e}{\partial\theta} \dfrac{\partial\theta}{\partial x_j} q_j \\
			& \qquad + D_{ij} \left[ \left( -\rho^2\left(\overline{\dfrac{\partial s}{\partial\rho}} 
			- \Lambda^e \dfrac{\partial e}{\partial\rho} \right) 
			- \Lambda^{\nabla\rho}_k \dfrac{\partial\rho}{\partial x_k} 
			+ \dfrac{\partial}{\partial x_k} \left( \Lambda^{\nabla\rho}_k\rho \right) \right) \delta_{ij} 
			% \right. \\ & \qquad \left. 
			+ \Lambda^e t_{ij} 
			- \Lambda^{\nabla\rho}_j \dfrac{\partial\rho}{\partial x_i} \right] \quad = \Sigma \geq 0.
	\end{align*}
	Since $\dfrac{\partial}{\partial x_k}\left(\Lambda^{\nabla\rho}_k\rho\right)=\dfrac{\partial\Lambda^{\nabla\rho}_k}{\partial x_k}\rho+\Lambda^{\nabla\rho}_k\dfrac{\partial\rho}{\partial x_k}$, the last inequality becomes 
	\begin{align*}
		& \dfrac{\partial}{\partial x_j} \left( \varphi_j - \Lambda^eq_j 
			- \Lambda^{\nabla\rho}_j\rho D_{kl}\delta_{kl} \right) 
			+ \dfrac{\partial\Lambda^e}{\partial\theta} \dfrac{\partial\theta}{\partial x_j} q_j 
		\\
		& \quad + D_{ij} \left[ \left( -\rho^2 \left( \overline{\dfrac{\partial s}{\partial\rho}} 
			- \Lambda^e \dfrac{\partial e}{\partial\rho}\right) + \dfrac{\partial\Lambda^{\nabla\rho}_k}{\partial x_k}\rho \right) \delta_{ij} 
			+ \Lambda^e t_{ij} - \Lambda^{\nabla\rho}_j \dfrac{\partial\rho}{\partial x_i} \right] 
			= \Sigma \geq 0,
	\end{align*}
	which proves \eqref{rezidualna nejednakost final Kortevegovi fluidi}.
\end{proof}

\end{document}
